\chapter{资料来源与版权说明}\label{app:sources}

本章（附录）汇总本书的资料来源构成，并说明本书在版权上的位置。
正文引用以 \ref{app:timeline} 页大事年表与参考文献列表为准，
此处只作\keyw{来源类型学}与\keyw{权利边界}的说明，不复述具体事实。

\section{来源类型与可信度}\label{app:sources:types}

本书的全部材料按产生方式分为四类，与正文的 A/B/C/M 四级对应：

\begin{table}[htbp]\centering
  \caption{材料来源类型与可信度对应}\label{tab:04:source}
  \small
  \begin{tabularx}{\textwidth}{@{}llX@{}}
    \toprule
    类型 & 可信度 & 说明 \\ \midrule
    一手权威报道 & \conf{A} & 如《人民日报》关于奶龙 IP 创新链的报道；
      可作为事实陈述直接引用 \\
    机构/平台公开信息 & \conf{B} & 如中国日报网转载的六周年活动通稿、
      百科与平台条目；可作事实陈述，注明来源与日期，
      多来源冲突时并列 \\
    社群考据与二手评述 & \conf{C} & 如失传媒体中文维基词条、
      自媒体考据文章；只能作为“线索”与“说法”，
      不得作为论证前提，正文须以“据……”“有说法认为”引出 \\
    本书自建模型 & \conf{M} & 本书为教学目的构造的框架、方程、
      示例数据与虚构问卷；必须显式声明非经验事实，
      参数取值为示意值 \\ \bottomrule
  \end{tabularx}
\end{table}

\begin{nlmethod}[使用纪律]
  \textbf{C 级材料的三条纪律}：
  (1) 不作为唯一依据支撑任何 A/B 级口吻的结论；
  (2) 涉及具体个人时，不转述其健康、生死、私德等个人境况的推测；
  (3) 当 C 级考据自身声明“来源缺失/仅为推测”时，
  本书照实保留这一自我限定，不替考据者“补全证据”。
\end{nlmethod}

\section{奶蛙起源考据的来源特别说明}\label{app:sources:naiwa}

奶蛙起源线索（第三章、第四章）主要依赖社群考据（\conf{C}）。
本书在引用时保留考据材料自身的三处自我限定：

\begin{enumerate}[label=\arabic*.]
  \item 考据指向的智能体（ID \fname{zNCvQDKfXy}）现已无法搜索、
        分享，评论区关闭，\keyw{不可复核}；
  \item 失传媒体维基词条自陈“待改进”“来源缺失”，
        且“星野 App 出现 bug”的结论“仅为推测、缺乏直接证据”；
  \item 多个生成路径推测均被考据者证伪，复现亦未成功。
\end{enumerate}

因此本书将奶蛙起源定性为一条\keyw{“未结案的证据链”}，
作为教学示范材料，而非已证事实。这一处理是全书
证据分级制度最严格的一次执行，读者可将其视为
“C 级材料应如何被使用”的标准样例。

\section{版权与合理使用}\label{app:sources:copyright}

\subsection{本书对形象的使用}

本书\keyw{不嵌入}受版权保护的原始动画截图、表情包原图、
周边商品图片。全部插图均为本书以 TikZ 自制的
抽象化示意图，不复制任何受保护的具体画面。
涉及“奶龙”动画形象、剧集名称及其相关商标时，
本书将其知识产权归属第七印象文化传媒（深圳）有限公司；
涉及“奶蛙”二次创作形象时，其署名权归属原始创作者。

\begin{nlcritique}[权利边界]
  本书对上述形象的使用限于\textbf{学术评论、教学示例与文化研究}
  目的，不主张任何权利，亦不构成对原作品著作权的替代或商业利用。
  正文将奶蛙二创的“未经授权形象再生产”作为\keyw{分析对象}
  加以讨论，不构成对任何具体侵权行为的背书或鼓励。
  如权利人认为本书某一处使用不当，请联系编写组予以替换或删除。
\end{nlcritique}

\subsection{引用社区用语与剧集名称}

本书引用的社群用语（如“糖”“倒放大法”“鲱鱼罐头”等）
与剧集名称（如《奶龙搞怪大作战》《奶龙之奇妙冒险》），
属于对公开传播内容的\keyw{指称性使用}，
用于识别与讨论特定文化对象，不构成对表达本身的复制。

\subsection{本书的再传播许可}

本书遵循\keyw{署名—非商业性使用—相同方式共享
（CC BY-NC-SA 4.0）}协议发布。各高校同好会以本校
活动材料为基础进行修订与续编时，须满足：
(1) 署名本书编写组与原出处；(2) 不用于商业目的；
(3) 衍生版本采用相同协议。修订时请同步更新
“事实手册”中相应条目的可信度标注与访问日期。

\section{访问日期与时效}

网络来源均记录首次访问日期（见参考文献列表）。
鉴于网络材料的时效性（“引文腐烂”问题，第十二章讨论），
读者在引用 \conf{B} / \conf{C} 级来源时，
应以本书记录的访问日期为基线，并在再引用时
重新核验链接的当前可达性与内容的现时状态。

\noindent\srcnote{本附录为版权与来源的“元说明”，
不替代正文的事实引用；具体引文格式以参考文献列表
（GB/T 7714—2015）为准。}
